\documentclass[aspectratio=169]{beamer}

% --- Tema e pacotes ---
\usetheme{Madrid}
\usecolortheme{seahorse}
\usepackage[utf8]{inputenc}
\usepackage[T1]{fontenc}
\usepackage{amsmath,amssymb}
\usepackage{graphicx}
\usepackage{booktabs}

% --- Metadados ---
\title[Criptografia Pós-Quântica]{Criptografia Pós-Quântica}
\author{Guilherme Moreira e Leonardo Martins}
\date{\today}

\begin{document}
	
	% --- Slide de título ---
	\begin{frame}
		\titlepage
	\end{frame}
	
	% --- O que é PQC ---
	\begin{frame}{O que é Criptografia Pós-Quântica (PQC)?}
		\begin{itemize}
			\item Algoritmos criptográficos seguros contra atacantes quânticos
			\item Executáveis em computadores clássicos
			\item Baseados em problemas matemáticos \emph{não} vulneráveis a Shor
			\item Foco principal: \textbf{criptografia de chave pública}
		\end{itemize}
	\end{frame}
	
		% --- Motivação ---
	\begin{frame}{Motivação}
		\begin{itemize}
			\item Algoritmos clássicos de chave pública (RSA, ECC)
			\item Ameaça de computadores quânticos escaláveis
			\item Algoritmos quânticos relevantes:
			\begin{itemize}
				\item \textbf{Shor}: fatoração e logaritmo discreto em tempo polinomial
				\item \textbf{Grover}: aceleração quadrática de busca exaustiva
			\end{itemize}
			\item Necessidade de \textbf{criptografia resistente a ataques quânticos}
		\end{itemize}
	\end{frame}
	
	
		\begin{frame}{Algoritmo de Shor}
		
		\begin{itemize}
			
			\item Algoritmo quântico proposto por \textbf{Peter Shor (1994)}
			
			\item Resolve em \textbf{tempo polinomial}:
			
			\begin{itemize}
				
				\item Fatoração de inteiros
				
				\item Logaritmo discreto
				
			\end{itemize}
			
			\item Base matemática:
			
			\begin{itemize}
				
				\item Redução do problema de fatoração ao \textbf{problema de ordem}
				
				\item Uso da \textbf{Transformada de Fourier Quântica (QFT)}
				
			\end{itemize}
			
		\end{itemize}
		
	\end{frame}
	
	
	% --- Impacto do Algoritmo de Shor ---
	
	\begin{frame}{Impacto do Algoritmo de Shor}
		
		\begin{itemize}
			
			\item Sistemas afetados:
			
			\begin{itemize}
				
				\item RSA (fatoração de \(N = pq\))
				
				\item Diffie-Hellman
				
				\item Criptografia de curvas elípticas (ECC)
				
			\end{itemize}
			
			\item Consequência:
			
			\begin{itemize}
				
				\item Chaves grandes \textbf{não aumentam a segurança}
				
				\item Segurança colapsa com computadores quânticos escaláveis
				
			\end{itemize}
			
			\item Motivação direta para a \textbf{Criptografia Pós-Quântica}
			
		\end{itemize}
		
	\end{frame} 
	
	% --- Algoritmo de Grover ---
	\begin{frame}{Algoritmo de Grover}
		\begin{itemize}
			\item Algoritmo quântico proposto por \textbf{ Lov Grover (1996)}
			\item \textbf{Aceleração Quadrática}: Otimiza buscas de $O(N)$ para $O(\sqrt{N})$.
			\item \textbf{Mecânica do Operador de Grover}:
			\begin{itemize}
				\item \textbf{Oráculo}: Identifica e inverte a fase do estado "marcado" (solução).
				\item \textbf{Difusão}: Realiza a "inversão sobre a média", amplificando a probabilidade da solução.
				\item \textbf{Iteração}: Repetido aproximadamente $\frac{\pi}{4}\sqrt{N}$ vezes para máxima probabilidade.
			\end{itemize}
			\item \textbf{Solução Pós-Quântica}: Dobrar o tamanho das chaves (ex: migrar para AES-256).
		\end{itemize}
	\end{frame}
	% --- Impacto do Grover ---
	\begin{frame}{Impacto do Algoritmo de Grover}
		\begin{itemize}
			\item Reduz segurança efetiva de chaves simétricas pela metade
			\item Exemplo:
			\begin{itemize}
				\item AES-128 \(\rightarrow\) ~64 bits de segurança quântica
				\item AES-256 recomendado no cenário pós-quântico
			\end{itemize}
			\item Hashes também sofrem impacto semelhante
			\begin{itemize}
				\item \textbf{SHA-256}: Resistência  a colisão cai para 128 
			\end{itemize}
		\end{itemize}
	\end{frame}
	 
	% --- Famílias de PQC ---
	\begin{frame}{Principais Famílias de PQC}
		\begin{enumerate}
			\item Criptografia baseada em \textbf{reticulados}
			\item Criptografia baseada em \textbf{códigos corretores de erro}
			\item Criptografia baseada em \textbf{funções hash}
		\end{enumerate}
	\end{frame}
	
	% --- Lattices ---
	\begin{frame}{Criptografia Baseada em Reticulados}
		\begin{itemize}
			\item Reticulado: conjunto discreto de pontos em \(\mathbb{R}^n\)
			\item Problemas difíceis:
			\begin{itemize}
				\item SVP – Shortest Vector Problem
				\item CVP – Closest Vector Problem
				\item LWE / Ring-LWE / Module-LWE
			\end{itemize}
			\item Vantagens:
			\begin{itemize}
				\item Boas provas de segurança
				\item Eficiência razoável
				\item Base da maioria dos padrões atuais
			\end{itemize}
		\end{itemize}
	\end{frame}
	
	% --- Exemplos Lattice ---
	\begin{frame}{Exemplos Baseados em Reticulados}
		\begin{itemize}
			\item Troca de chaves / KEM:
			\begin{itemize}
				\item Kyber
			\end{itemize}
			\item Assinaturas digitais:
			\begin{itemize}
				\item Dilithium
				\item Falcon
			\end{itemize}
		\end{itemize}
	\end{frame}
	
	% --- Code-based ---
	\begin{frame}{Criptografia Baseada em Códigos}
		\begin{itemize}
			\item Baseada em códigos corretores de erro
			\item Problema difícil: \textbf{decodificação de códigos lineares aleatórios}
			\item Algoritmo clássico:
			\begin{itemize}
				\item McEliece
			\end{itemize}
			\item Características:
			\begin{itemize}
				\item Segurança estudada há décadas
				\item Chaves públicas muito grandes
			\end{itemize}
		\end{itemize}
	\end{frame}
	
	% --- Hash-based ---
	\begin{frame}{Criptografia Baseada em Hash}
		\begin{itemize}
			\item Segurança baseada apenas em funções hash
			\item Principal uso: \textbf{assinaturas digitais}
			\item Exemplos:
			\begin{itemize}
				\item XMSS
				\item SPHINCS+
			\end{itemize}
			\item Propriedades:
			\begin{itemize}
				\item Segurança conservadora
				\item Assinaturas maiores
				\item Algumas variantes são stateful
			\end{itemize}
		\end{itemize}
	\end{frame}
	
	% --- Padronização ---
	\begin{frame}{Padronização e NIST}
		\begin{itemize}
			\item Processo de padronização do NIST (EUA)
			\item Objetivo: substituir RSA e ECC
			\item Seleções principais:
			\begin{itemize}
				\item Kyber (KEM)
				\item Dilithium, Falcon, SPHINCS+ (assinaturas)
			\end{itemize}
		\end{itemize}
	\end{frame}
	
	% --- Desafios ---
	\begin{frame}{Desafios da PQC}
		\begin{itemize}
			\item Tamanhos de chaves e assinaturas
			\item Integração com protocolos existentes (TLS, PKI)
			\item Implementações seguras
			\item Resistência a ataques de canal lateral
		\end{itemize}
	\end{frame}
	
	% --- Conclusão ---
	\begin{frame}{Conclusão}
		\begin{itemize}
			\item Computação quântica muda o cenário criptográfico
			\item PQC é essencial para segurança de longo prazo
			\item Reticulados são a abordagem dominante
			\item Migração deve começar antes da ameaça prática
		\end{itemize}
	\end{frame}
	
	% --- Slide final ---
	\begin{frame}
		\centering
		\Huge Perguntas?
	\end{frame}
	
\end{document}
